\section{Resident Performance Under Explicit Assumptions}
\subsection{A Separate Optimized Software Cohort}
Resident completion includes state initialization, cost loads, all waves, canonical/directory retirement, output writes and acknowledgment drain. Host preparation and full readback are separate. This boundary differs from native E2 and Process. We reuse an affinity/NUMA-bound packed CPU implementation on dual Xeon Gold 6530 processors, testing 8, 32 and 64 workers. The 64-worker setting minimizes mean resident time on both full 40-sample workloads. The allocation exposes 64 physical cores with a 55-CPU-time quota; 40 samples allow at most 40 simultaneous solves. This is the strongest measured CPU in this allocation, not a claim about every CPU implementation.

Three independent processes each execute one warmup and three formal repetitions. Table~\ref{tab:resident} reports the mean and sample SD of the three process means, not nine independent trials. The retained H100 implementation exposes sample parallelism and actual 64-bit atomic whole-pair minima; its resident interval includes initialization, launch/control, completion reads and final synchronization. Full output transfer and research comparison are outside resident. The older BNN GPU receipt lacks a contemporaneous executable digest; retained source/build records and the current executable digest are not misrepresented as an execution-time hash.

\begin{table}[t]\centering
\caption{Complete-batch resident software time (seconds). CPU dispersion is over three independent process means. This cohort is separate from native eight-core E2/Process.}\label{tab:resident}
\setlength{\tabcolsep}{4pt}
\begin{tabular}{lrr}\toprule & BNN40 & LSTM40\\\midrule
CPU (64 workers) & $0.699624\pm0.026647$ & $0.725510\pm0.021287$\\
Retained H100 mean & 0.151800 & 0.159498\\
\bottomrule\end{tabular}
\end{table}

\subsection{Finite Model: Assumed Frequency and Memory}
We evaluate a \emph{hypothetical 200 MHz core} with declared 32-byte memory service. Baseline latency/issue interval is 80/5 ns; a conservative sensitivity scenario uses 160/10 ns. These correspond to 16/1 and 32/2 cycles at the assumed frequency. Reads and writes share the 32 physical ports; finite queues, credits, return identity and blocking remain active. The ideal issue ceilings of 204.8 and 102.4 GB/s are not measured random-RMW HBM throughput, and the conservative scenario is not a worst-case guarantee.

\begin{table}[t]\centering
\caption{Complete finite-model results at \emph{assumed 200 MHz}; neither routed-frequency performance nor board measurement.}\label{tab:model}
\begin{tabular}{llrr}\toprule Case & Latency/issue & Cycles & Model s\\\midrule
BNN40 & 80/5 ns & 102,557,024 & 0.512785\\
LSTM40 & 80/5 ns & 104,154,414 & 0.520772\\
BNN40 & 160/10 ns & 135,128,351 & 0.675642\\
LSTM40 & 160/10 ns & 137,198,248 & 0.685991\\
\bottomrule\end{tabular}
\end{table}

These deterministic predictions have no invented sampling error bars. Baseline expansion consumes 77.09/77.15\% of BNN/LSTM cycles and finalization 20.95/20.89\%; all other mutually exclusive phases account for the remainder. Candidate throughput is approximately 16.28/16.30 per expansion cycle, about 25.4\% of the 64-lane ideal. Finite scheduling, transport and publication, together with sequential FPGA sample execution, explain why lane count alone is insufficient. Overlapping stall counters cannot be added as independently removable time.

The conservative BNN prediction exceeds the fastest measured CPU process mean, 0.668921 s, despite lying below its three-process mean. Thus even the hypothetical-frequency result does not establish a robust CPU advantage. The GPU remains faster in both assumed scenarios; beating it is not a qualification requirement. Frequency changes alter service discretization and scheduling: $T=C(f,\text{service},\text{structure})/f$, with absolute nanoseconds rounded up to cycles. Dividing old cycles by a new frequency is invalid when that binding changes.

Resident predictions are not native E2 estimates. Required packed-descriptor construction is shared by the selected CPU and FPGA representations. Host output-layout conversion was measured separately at 1.561/1.537 ms on preallocated, prefaulted buffers; it is not board DMA and excludes transfers and native map assembly. Such costs cannot be charged only to FPGA or combined with unrelated native timing cohorts to fabricate an end-to-end speedup.
